% Generated from the matching Typst scaffold by
% scripts/migrate_typst_report_to_latex.py (prose drafted 2026-06-25).

% ── 1.1  Background and Motivation ──────────────────────────────────────────
% PURPOSE: Motivate persistent memory for LLM coding agents on long-horizon work.
% MOVE: Scale+stakes -> traditional approach -> structural limitation.
% ─────────────────────────────────────────────────────────────────────────

\section{Background and Motivation}

Large language model (LLM) coding agents have entered the SWE-Bench era: given a
natural-language issue and a real repository, an agent now plans, edits files, runs the
test suite, and iterates until the issue is resolved, autonomously and end to
end~\cite{joshi2025}. The capability is no longer a laboratory curiosity --- agents
resolve a non-trivial fraction of genuine GitHub issues drawn from mature open-source
projects --- and the stakes scale with it: the same automation that closes one issue is
expected to work an entire backlog.

Real software-engineering work, however, is rarely a sequence of isolated one-shot
prompts. It is long-horizon and \emph{sequential}: many related tasks accrue over the
lifetime of a single codebase, each leaving behind a diagnosis, a fix, and a lesson that
might inform the next. An agent that treats every issue as if it were the first throws
away precisely the experience a human maintainer would carry forward. The natural remedy
is to externalise that experience as \emph{persistent memory} --- a store of past
problems and their solutions that the agent can consult on future
tasks~\cite{packer2024,zhong2024}. This is the premise behind a growing family of agent
memory systems, and it is intuitively compelling: an agent that remembers more of its
past should solve more of its future.

The traditional way to operationalise that premise is simply to keep everything. Each
completed task appends a record; nothing is ever discarded. But naive accumulation has a
structural cost. Context windows are finite and, more importantly, they are not free to
fill: longer prompts raise token cost and latency, and a well-documented body of evidence
shows that model performance degrades as relevant information is buried among
distractors~\cite{liu2024} and even as raw context length grows with retrieval held
perfect~\cite{du2025}. Worse, more memory is not merely more expensive --- it can be
actively harmful. Recent empirical work shows that agents exhibit
\emph{experience-following}, faithfully reusing retrieved precedents even when those
precedents mislead, so that a larger or lower-quality memory store degrades
behaviour rather than improving it~\cite{xiong2025}.

This sets up the structural limitation that motivates the thesis. If keeping everything is
costly and sometimes harmful, then an agent ought to forget --- but \emph{forgetting} is
itself a policy decision with no obvious right answer, and it is entangled with how memory
is retrieved. Does retained experience actually help downstream coding performance, or does
it mostly inflate the prompt? And if some forgetting is warranted, \emph{when} and
\emph{how much}? These questions have been asked for conversational and
general-purpose agents, but they have not been answered cleanly for coding agents under a
controlled comparison. Figure~\ref{fig:architecture} situates the question: storage and
forgetting are one component in a pipeline that also includes retrieval, and only by
holding the rest fixed can the effect of forgetting be read off in isolation.
